\documentclass[11pt]{article}

\usepackage[T1]{fontenc}
\usepackage[utf8]{inputenc}
\usepackage{lmodern}
\usepackage[margin=1.1in]{geometry}
\usepackage{amsmath,amssymb,amsthm,amscd,mathtools}
\usepackage{mathrsfs}
\usepackage{microtype}
\usepackage{xcolor}
\usepackage{array}
\usepackage{longtable}
\usepackage{enumitem}
\usepackage{titlesec}
\usepackage{hyperref}
\usepackage{xurl}

\titleformat{\section}
  {\normalfont\LARGE\bfseries}{\thesection}{1em}{}
\titleformat{\subsection}
  {\normalfont\Large\bfseries}{\thesubsection}{1em}{}

\hypersetup{
  pdftitle={A Formal Projector Spine for the Rational Hodge Target},
  pdfauthor={Chloe van der Vlugt},
  colorlinks=true,
  linkcolor=blue!55!black,
  citecolor=green!40!black,
  urlcolor=blue!60!black
}

\newtheorem{theorem}{Theorem}[section]
\newtheorem{proposition}[theorem]{Proposition}
\newtheorem{corollary}[theorem]{Corollary}
\theoremstyle{definition}
\newtheorem{definition}[theorem]{Definition}
\theoremstyle{remark}
\newtheorem*{remark}{Remark}

\DeclareMathOperator{\Spec}{Spec}
\DeclareMathOperator{\Frac}{Frac}
\DeclareMathOperator{\Gal}{Gal}
\DeclareMathOperator{\Corr}{Corr}
\DeclareMathOperator{\End}{End}
\DeclareMathOperator{\ev}{ev}
\DeclareMathOperator{\CH}{CH}
\DeclareMathOperator{\Ram}{Ram}
\DeclareMathOperator{\Sat}{Sat}
\DeclareMathOperator{\Ext}{Ext}
\DeclareMathOperator{\Cl}{Cl}
\DeclareMathOperator{\Lift}{Lift}
\DeclareMathOperator{\Def}{Def}

\newcommand{\Q}{\mathbf Q}
\newcommand{\one}{\mathbf 1}
\newcommand{\fe}{\mathrm{f\acute et}}
\newcommand{\Geta}{\Gamma_{\eta}}
\newcommand{\Gbar}{\overline{\Gamma}}
\newcommand{\formal}{\textnormal{\textsc{Formal}}}
\newcommand{\conditional}{\textnormal{\textsc{Conditional}}}
\newcommand{\openstatus}{\textnormal{\textsc{Open}}}
\newcommand{\sourceinput}{\textnormal{\textsc{Source Input}}}
\newcommand{\leanpath}[1]{\texttt{\nolinkurl{#1}}}
\newcommand{\corrsrc}[1]{%
  \leanpath{HodgeConjecture/Correspondence/#1}}
\newcommand{\collsrc}[1]{%
  \leanpath{HodgeConjecture/Collision/#1}}
\newcommand{\sheetsrc}[1]{%
  \leanpath{HodgeConjecture/Boundary/Sheets/#1}}
\newcommand{\normsrc}[1]{%
  \leanpath{HodgeConjecture/Boundary/Normalization/#1}}
\newcommand{\boundarysrc}[1]{%
  \leanpath{HodgeConjecture/Boundary/#1}}
\newcommand{\chowsrc}[1]{%
  \leanpath{HodgeConjecture/Realization/Chow/#1}}

\newenvironment{leanfiles}
  {\par\medskip\begingroup\footnotesize\raggedright\noindent
   \textit{Lean correspondence.}\ }
  {\par\endgroup\medskip}

\newenvironment{statusnote}[1]
  {\par\medskip\noindent #1.\ \begingroup}
  {\par\endgroup\medskip}

\setlength{\emergencystretch}{1.5em}

\title{A Formal Projector Spine for the Rational Hodge Target:\\
  Collisions, Marked Galois Sheets, and Boundary Interfaces}
\author{Chloe van der Vlugt}
\date{Draft of September 11, 2026}

\begin{document}

\maketitle

\begin{abstract}
We describe a Lean~4 development organized around the exact abstract target
of the rational Hodge conjecture.  For a supplied cycle-class model
\(\mathscr H=(A,V,V_{\mathrm{Hdg}},\operatorname{cl})\), the modeled target
\(\mathrm{HC}_{\Q}(\mathscr H)\) says that every class in the independently
specified Hodge subspace lies in the image of the cycle-class map.  The
development proves that, after supplying explicit Chow-level construction
and lower-left detection maps, this target is equivalent to lifting every
canonical graph projector \(P_\alpha(t,x)=(t,t\alpha)\).

The dependency-ordered components leading toward such lifts consist of a collision algebra
and its separability decomposition, a free finite-\'etale span algebra,
transport of group and Hecke idempotents, a marked \(G/H\)-sheet calculation
putting a selector into graph-block normal form, normal-closure deck actions,
and a conjugate-boundary separation theorem excluding ramification in
codimension one.  For compactification we distinguish a raw closure, whose
idempotence defect merely restricts to zero, from a supplied compactified
deck-graph square, which is a sufficient route to a genuinely idempotent
compactified image.

These statements are machine checked from their displayed hypotheses
(\formal).  The geometric comparison, compactified realization, and
class-dependent marked-sheet coverage packages remain supplied inputs
(\conditional), and constructing them for genuine Chow groups and Betti
cohomology remains \openstatus.  In particular, this manuscript proves no
new case, partial case, or unconditional instance of the Hodge conjecture.
\end{abstract}

\begin{center}
  \small Canonical manuscript candidate for author review.  No novelty,
  priority, authorship-finality, or license claim is made in this draft.
\end{center}

\begingroup
\small
\tableofcontents
\endgroup

\section{Target, status convention, and proof spine}
\label{sec:target}

The manuscript uses three theorem-status labels.  \formal\ means that a Lean theorem
checks from the hypotheses printed in its type.  \conditional\ means that
one or more mathematically substantive geometric objects or comparison
statements are supplied rather than constructed by the repository.
\openstatus\ marks the absence of the construction needed to instantiate an
interface in the intended Hodge-theoretic setting.  A result may therefore be
both \formal\ and \conditional: its implication is machine checked, while a
key premise is still an input.  A fourth annotation, \sourceinput, marks a
published theorem cited in place of repeating its proof; until it is encoded
in the dependency library, the Lean endpoint retains its consequence as a
named premise.

Fix rational vector spaces \(A\) and \(V\), a subspace
\(V_{\mathrm{Hdg}}\leq V\), and a linear cycle-class map

\[
  \operatorname{cl}:A\longrightarrow V,
  \qquad
  \operatorname{im}(\operatorname{cl})\leq V_{\mathrm{Hdg}}.
\]

The datum
\(\mathscr H=(A,V,V_{\mathrm{Hdg}},\operatorname{cl})\) is a
\emph{rational Hodge cycle-class model}.  Its Hodge subspace is independent
of the range of \(\operatorname{cl}\); thus the conjecture is not built into
the model.  Define

\begin{align*}
  \operatorname{Hdg}_{\Q}(\alpha)&:\Longleftrightarrow
    \alpha\in V_{\mathrm{Hdg}},\\
  \operatorname{Alg}_{\Q}(\alpha)&:\Longleftrightarrow
    \alpha\in\operatorname{im}(\operatorname{cl}),\\
  \mathrm{HC}_{\Q}(\mathscr H)&:\Longleftrightarrow
    V_{\mathrm{Hdg}}\leq\operatorname{im}(\operatorname{cl}).
\end{align*}

For a genuine smooth projective complex variety \(X\) and codimension \(p\),
the intended specialization is

\[
  A=\CH^p(X)_{\Q},\qquad
  V=H^{2p}(X,\Q),\qquad
  V_{\mathrm{Hdg}}=H^{2p}(X,\Q)\cap H^{p,p}(X),
\]

with \(\operatorname{cl}\) the rational cycle-class map.  The repository
does not yet construct these objects.  The family-level Lean statement
recovers the classical universal quantification only if its index type is
actually populated by all smooth projective complex varieties and its models
are the genuine geometric models.

Let \(\Gamma\) be a rational algebra acting on \(\Q\oplus V\) through
\(\rho:\Gamma\to\End_{\Q}(\Q\oplus V)\).  For \(\alpha\in V\), define

\[
  P_\alpha(t,x)=(t,t\alpha).
\]

Then \(P_\alpha^2=P_\alpha\), and the lower-left evaluation

\[
  d_{\mathrm{ll}}(T)
  =\operatorname{pr}_V\bigl(T(1,0)\bigr)
\]

recovers \(d_{\mathrm{ll}}(P_\alpha)=\alpha\).  Supply two independent
interfaces:

\begin{enumerate}[label=\textup{(\roman*)},leftmargin=2.6em]
  \item a construction \(C:A\to\Gamma\) such that \(C(z)\) is idempotent
        and \(\rho(C(z))=P_{\operatorname{cl}(z)}\);
  \item a linear detection map \(D:\Gamma\to A\) such that
        \(\operatorname{cl}(D(\gamma))=d_{\mathrm{ll}}(\rho(\gamma))\).
\end{enumerate}

Write \(\Lift_{\mathrm{CH}}(P_\alpha)\) for the type of idempotents
\(\Pi\in\Gamma\) whose action is \(P_\alpha\).

\begin{theorem}[Graph-projector equivalence --- \formal]
\label{thm:hodge-graph-equivalence}
With construction and detection as above,
\[
  \mathrm{HC}_{\Q}(\mathscr H)
  \quad\Longleftrightarrow\quad
  \forall\alpha\in V_{\mathrm{Hdg}},\quad
    \Lift_{\mathrm{CH}}(P_\alpha)\neq\varnothing.
\]
The equivalence is also formalized degreewise for one model and for an
indexed family of models.  If a separate predicate of Hodge morphisms is
supplied together with
\(P_\alpha\text{ Hodge}\Longleftrightarrow
\alpha\in V_{\mathrm{Hdg}}\), the right side can equivalently quantify over
Hodge graph projectors.
\end{theorem}

\begin{proof}
If \(\alpha=\operatorname{cl}(z)\), the idempotent \(C(z)\) acts as
\(P_\alpha\).  Conversely, if \(\Pi\) acts as \(P_\alpha\), then
\(\operatorname{cl}(D(\Pi))=d_{\mathrm{ll}}(P_\alpha)=\alpha\).
The quantified statement is obtained pointwise.
\end{proof}

\begin{leanfiles}
The model and target are \leanpath{RationalHodgeCycleClassModel} and
\leanpath{RationalHodgeConjectureAt} in
\leanpath{HodgeConjecture/Realization/HodgeClasses.lean}.  The graph
operator, construction and detection interfaces, and
Theorem~\ref{thm:hodge-graph-equivalence} are in
\chowsrc{TateGraphProjector.lean}, principally
\leanpath{ClassGraphProjector.operator},
\leanpath{ClassGraphProjectorConstruction},
\leanpath{ClassGraphProjectorDetection}, and
\leanpath{rationalHodgeConjectureAt_iff_everyHodgeClassHasGraphProjectorLift}.
The degreewise and family forms are named analogously.
\end{leanfiles}

\begin{statusnote}{\conditional/\openstatus}
Theorem~\ref{thm:hodge-graph-equivalence} is an exact logical reduction, not
an instance of the conjecture.  Construction, detection, genuine Chow and
cohomology objects, and the Hodge-morphism comparison remain to be built in
the intended geometry.
\end{statusnote}

The remaining sections present the active dependency-ordered components
toward the right side of Theorem~\ref{thm:hodge-graph-equivalence}.  They
include parallel compactification branches and do not constitute a proved
composite arrow from boundary separation to Hodge coverage.  In particular,
they do not replace the theorem's explicit construction and detection
hypotheses.

\section{Collision algebra, separability, and finite-\'etale spans}
\label{sec:collision}

For a commutative algebra \(B\to R\), its self-collision algebra and diagonal
map are

\[
  \mathcal A_{R/B}=R\otimes_B R,
  \qquad
  \mu:\mathcal A_{R/B}\longrightarrow R,
  \qquad x\otimes y\longmapsto xy.
\]

The diagonal obstruction is \(J_\Delta=\ker\mu\), equivalently the ideal
generated by the differences \(1\otimes r-r\otimes1\), and the residual
off-diagonal ideal is

\[
  J_{\mathrm{off}}=\operatorname{Ann}(J_\Delta).
\]

The map \(\mu\) is surjective by the section \(r\mapsto1\otimes r\), hence
\(\mathcal A_{R/B}/J_\Delta\simeq R\).  When \(B\to R\) is essentially of
finite type and formally unramified, the contact and overlap defects vanish;
on affine spectra this gives a decomposition of the actual self-fiber
product into the diagonal and residual off-diagonal pieces.

\begin{proposition}[Affine collision decomposition --- \formal]
\label{prop:collision-decomposition}
Under essential finite type and formal unramifiedness,
\[
  \Spec R\times_{\Spec B}\Spec R
  \simeq
  \Spec R\ \amalg\
  \Spec\bigl(\mathcal A_{R/B}/J_{\mathrm{off}}\bigr),
\]
and the first coproduct inclusion is the relative diagonal.
\end{proposition}

\begin{leanfiles}
The algebra, diagonal map, ideals, and surjectivity are in
\collsrc{Algebra.lean}.  The scheme comparison and formal-unramified
decomposition are in \collsrc{AffineGeometry.lean}, under
\leanpath{collisionDiagonalQuotientEquiv},
\leanpath{affineCollisionDefects_vanish_of_formallyUnramified}, and
\leanpath{affineSelfCollisionDecompositionOfFormallyUnramified}.
\end{leanfiles}

For a morphism \(p_X:X\to T\), an endocorrespondence is a span
\(X\xleftarrow{\ell}Z\xrightarrow{r}X\) over \(T\).  Fiber-product
composition is associative and unital up to canonical isomorphism; quotienting
by carrier isomorphism makes these laws strict.  Requiring both legs to be
finite \'etale is stable under identity, transpose, base change, and
composition.  Thus the admissible classes form a monoid, and its rational
monoid algebra is

\[
  \mathcal C^{\mathrm{free}}_{\fe}(X/T)
  =\Q[\Corr_{\fe}(X/T)].
\]

\begin{theorem}[Free finite-\'etale convolution algebra --- \formal]
\label{thm:free-correspondence}
The algebra \(\mathcal C^{\mathrm{free}}_{\fe}(X/T)\) is unital and is not
assumed commutative.  Transpose is an involution satisfying
\((ab)^{\mathsf t}=b^{\mathsf t}a^{\mathsf t}\).  A separately supplied
\leanpath{FiniteEtaleAlgebraModel} can impose finite additivity and map this
free algebra to another rational algebra.
\end{theorem}

\begin{leanfiles}
Raw spans, classes, admissibility, and transpose occur in
\corrsrc{Scheme.lean}, \corrsrc{Classes.lean},
\corrsrc{EtaleClasses.lean}, and \corrsrc{TransposeAlgebra.lean}.
The additive realization interface is in \corrsrc{FiniteEtaleModel.lean}.
\end{leanfiles}

\begin{statusnote}{\formal/\conditional}
The collision and free-span constructions are machine checked.  The free
span algebra is not a Chow correspondence algebra: it imposes neither
rational equivalence nor a relation identifying a finite coproduct with a
sum.  Such additivity is intentionally a separate model field.
\end{statusnote}

\section{Group, collision, and Hecke projectors}
\label{sec:group}

An action of a group \(G\) on \(X\) over \(T\) sends \(g\) to its graph.
Graph composition gives an algebra homomorphism

\[
  \gamma:\Q[G]\longrightarrow
  \mathcal C^{\mathrm{free}}_{\fe}(X/T),
  \qquad [g]\longmapsto[\Gamma_g].
\]

\begin{theorem}[Transport of group projectors --- \formal]
\label{thm:group-transport}
Every idempotent \(e\in\Q[G]\) maps to a convolution idempotent.  Finite
ordered orthogonal decompositions of \(1\) map to decompositions of the
diagonal.  If \(G\) is finite, the average
\[
  e_G=\frac1{|G|}\sum_{g\in G}[g]
\]
and \(1-e_G\) map to complementary projectors.
\end{theorem}

For a finite subgroup \(H\leq G\), put

\[
  e_H=\frac1{|H|}\sum_{h\in H}[h],
  \qquad
  \mathcal H(G,H)=e_H\Q[G]e_H.
\]

The corner has unit \(e_H\), and the normalized sandwiched element
\(T_g=e_H[g]e_H\) depends only on \(HgH\).

\begin{theorem}[Hecke-corner transport --- \formal]
\label{thm:hecke-transport}
The graph algebra map restricts to a unital ring homomorphism from
\(\mathcal H(G,H)\) to the corresponding corner of the finite-\'etale
correspondence algebra.  Hence every supplied Hecke-corner idempotent gives a
formal convolution projector.
\end{theorem}

If a finite-\'etale cover is supplied with a \(G\)-torsor collision
decomposition

\[
  X\times_YX\simeq\coprod_{g\in G}X
\]

compatible with the graph on each summand, then any supplied finite-additive
model \(\chi\) satisfies

\[
  \chi([X\times_YX])=\sum_{g\in G}\chi([\Gamma_g]).
\]

Thus \(|G|^{-1}\chi([X\times_YX])\) is the image of \(e_G\) and is
idempotent.

\begin{statusnote}{\formal/\conditional}
The transport theorems are formal.  The Galois collision formula is formal
from a supplied torsor decomposition and finite-additive realization; those
data are not inferred merely from the label ``finite \'etale.''
\end{statusnote}

\begin{leanfiles}
The graph map and group-idempotent transport are in
\corrsrc{Graph.lean}, \corrsrc{FiniteGroupProjector.lean}, and
\corrsrc{GroupAlgebraProjector.lean}.  The Hecke construction is in
\corrsrc{HeckeProjector.lean}, particularly
\leanpath{SubgroupHeckeAlgebra}, \leanpath{heckeElement_eq_of_rel}, and
\leanpath{heckeCorrespondence_isIdempotent}.  The conditional collision
comparison is in \corrsrc{GaloisCollision.lean}.
\end{leanfiles}

\section{Marked sheets and the block graph normal form}
\label{sec:block}

Let \(G\) be a group and \(H\leq G\).  Keep the pointwise sheet algebra

\[
  F_k(H)=k^{G/H}
\]

distinct from \(\Q[G]\) and from any compactified correspondence algebra.
Let \(q_H\) be the characteristic function of the base coset, and set

\[
  E_H=\ker(\ev_H:F_{\Q}(H)\to\Q),
  \qquad a_H=q_H-\one\in E_H.
\]

There are two splittings of base-sheet evaluation.  The selector splitting
uses \(q_H\), while the constant-unit splitting uses \(\one\).  In formulas,

\begin{align*}
  \Phi_{q_H}(f)&=(f(H),f-f(H)q_H),\\
  \Phi_1(f)&=(f(H),f-f(H)\one).
\end{align*}

Let \(m_{q_H}\) denote pointwise multiplication by \(q_H\), and define

\[
  P_0(t,x)=(t,0),\qquad
  P_{a_H}(t,x)=(t,ta_H),\qquad
  U_{a_H}(t,x)=(t,x+ta_H).
\]

\begin{theorem}[Marked-sheet block normal form --- \formal]
\label{thm:block}
The following identities hold without a finiteness assumption on \(G/H\):
\begin{align*}
  \Phi_{q_H}m_{q_H}\Phi_{q_H}^{-1}&=P_0,\\
  \Phi_1m_{q_H}\Phi_1^{-1}&=P_{a_H},\\
  \Phi_1\Phi_{q_H}^{-1}&=U_{a_H},\\
  P_{a_H}&=U_{a_H}P_0U_{a_H}^{-1}.
\end{align*}
In particular, selector multiplication becomes the same abstract graph
operator used in Theorem~\ref{thm:hodge-graph-equivalence}.
\end{theorem}

For any linear map \(\theta:E_H\to V\), put
\(\Theta=\operatorname{id}_{\Q}\oplus\theta\) and
\(\alpha=\theta(a_H)\).  Naturality gives

\[
  \Theta P_{a_H}=P_\alpha\Theta.
\]

This identity creates the precise linear bridge from a marked selector to a
class graph.  It does not construct \(\theta\), make it motivic, or show that
\(\alpha\) is algebraic.

\begin{leanfiles}
The ring-generic graph operator, diagonal operator, shear, and conjugacy
identity are shared through
\leanpath{HodgeConjecture/LinearAlgebra/BlockGraphProjector.lean}.  The sheet
algebra and selector are in \sheetsrc{GenericProjector.lean}.  The two
specialized splittings and selector-multiplication identities are in
\sheetsrc{MarkedBlockProjector.lean}.  Their identification with the Tate
graph operator and naturality under \(\theta\) are in
\chowsrc{GaloisMarkedBlock.lean}, under
\leanpath{graphOperator_eq_classGraphProjector} and
\leanpath{productRealization_intertwines_selectorMultiplication}.
\end{leanfiles}

\section{Normal closure, conjugate boundaries, and codimension one}
\label{sec:boundary}

Let \(B\) be a domain with \(K=\Frac(B)\), let \(L/K\) be finite, and
supply a normal closure \(N/K\) with a marked embedding \(\iota:L\hookrightarrow
N\).  Write

\[
  G=\Gal(N/K),\qquad H=\Gal(N/\iota(L)),
\]

and use the canonical normalization notation

\[
  \widetilde B_L=\operatorname{IntCl}_B(L),\quad
  \widetilde B_N=\operatorname{IntCl}_B(N),\qquad
  Y_L=\Spec(\widetilde B_L),\quad
  Y_N=\Spec(\widetilde B_N).
\]

The group \(G\) acts on \(\widetilde B_N\) and hence on \(Y_N\) over
\(\Spec B\).  The fixed normalization is canonically

\[
  \operatorname{IntCl}_B(N^J)
  \simeq_B(\widetilde B_N)^J
\]

for every subgroup \(J\leq G\); under the marked perfect-field hypotheses,
\(\widetilde B_L\simeq_B(\widetilde B_N)^H\).

Each \(g\in G\) gives a conjugate embedding \(\iota_g:L\hookrightarrow N\)
and a canonical leg \(c_g:Y_N\to Y_L\), depending only on \(gH\).  For a
marked normalized cover
\(\mathcal C=(U\hookrightarrow Y_L,\mathcal N)\), let

\[
  \partial U=Y_L\setminus U,\qquad
  \partial_g=c_g^{-1}(\partial U),\qquad
  J_g=I(\partial_g).
\]

For \(I\leq G\), define

\begin{align*}
  J_{\mathrm{fix}}(I)
    &=(\sigma(t)-t)_{\sigma\in I,\,t\in\widetilde B_N},\\
  J_{\mathrm{mov}}(I)
    &=\sum_{I\nleq gHg^{-1}}J_g,\\
  J_{\mathrm{fm}}(I)
    &=J_{\mathrm{fix}}(I)+J_{\mathrm{mov}}(I).
\end{align*}

Conjugate-boundary separation is the condition

\[
  \operatorname{Sep}_{\partial}(\mathcal C):
  \quad \forall I\leq G,\ \forall Q\in\Spec(\widetilde B_N),\quad
  I\neq1,\ \operatorname{ht}(Q)=1
  \Longrightarrow J_{\mathrm{fm}}(I)\nsubseteq Q.
\]

Let \(\mathscr D\) be the normalization/inertia diagram, let
\(\operatorname{Ctr}(\mathscr D)\) identify its centers with the canonical
legs \(c_g\), and let \(\operatorname{RamCmp}(\mathscr D)\) identify its
relative inertia indices with the geometric ramification indices.  These are
independent comparisons, bundled as \(\mathscr B_{\mathscr D}\).

\begin{theorem}[Conjugate-boundary endpoint --- \formal; \conditional]
\label{thm:boundary-endpoint}
Assume that \(B\) is Noetherian, that
\(\mathscr B_{\mathscr D}\) is supplied, that the marked visible morphism
\(U\to\Spec B\) is formally unramified, and that
\(\operatorname{Sep}_{\partial}(\mathcal C)\) holds.  Then
\[
  \operatorname{Ram}^{(1)}(Y_N/B)=\varnothing;
\]
equivalently, every height-one point of \(Y_N\) is unramified over \(B\).
\end{theorem}

\begin{proof}
A ramified height-one point has nontrivial valuation inertia \(I\).
Residue-field triviality puts \(J_{\mathrm{fix}}(I)\) in its prime.  For
every sheet moved by \(I\), the center and ramification comparisons plus
formal unramifiedness put the corresponding \(J_g\) in the same prime.
Thus \(J_{\mathrm{fm}}(I)\) is contained in that height-one prime, contrary
to separation.
\end{proof}

\begin{leanfiles}
The marked normal closure is in \normsrc{NormalClosure.lean}; its action and
fixed-ring comparison are in \normsrc{DeckAction.lean} and
\normsrc{InvariantIntermediate.lean}.  The conjugate maps are in
\normsrc{ConjugateNormalization.lean}.  The ideals, separation predicate,
and endpoint are in \normsrc{ConjugateBoundary.lean}, chiefly
\leanpath{ConjugateBoundarySeparation} and
\leanpath{noCodimensionOneRamification_of_conjugateBoundarySeparation}.
The empty-locus and proposition presentations are connected by explicit
conversion theorems rather than treated as separate geometric conclusions.
\end{leanfiles}

\begin{statusnote}{\openstatus}
The repository does not prove the separation hypothesis in the intended
compactification, and the codimension-one conclusion is not itself a theorem
that a cover, selector, or Chow correspondence extends globally.  Any such
purity or gluing arrow remains an explicit additional input.
\end{statusnote}

\begin{statusnote}{\sourceinput/\openstatus}
A precise standard cover-extension theorem is available without repeating
its proof.  If \(X\) is locally Noetherian, \(U\subset X\) is dense open,
\(Y\to U\) is finite \'etale and unramified over \(X\) in codimension one,
and every local ring of \(X\setminus U\) is regular, Stacks Project
Lemma~58.31.2 produces a finite \'etale extension over \(X\)
\cite{StacksFiniteEtaleCodimOneExtension}.  Applying it here still requires
an exact identification of \(X,U,Y\) and proofs of its regularity and
codimension-one hypotheses.  If, in addition, the underlying topological
space of \(X\) is Noetherian and, for every \(x\in X\setminus U\), the
punctured spectrum of the strict henselization of \(\mathcal O_{X,x}\) is
connected, restriction of finite \'etale covers to the dense open is fully
faithful \cite{StacksFiniteEtaleFullFaithful}; hence
deck automorphisms of a chosen extension extend uniquely.  Neither citation
constructs the compactified Chow correspondence or its realization square.
\end{statusnote}

The generic normal-closure evaluation is nevertheless formal.  The cosets
\(G/H\) index the marked embeddings and give

\[
  \epsilon:L\otimes_KL\longrightarrow (G/H\to N),
  \qquad
  \epsilon(x\otimes y)(gH)=\iota(x)g(\iota(y)).
\]

For a formally unramified, essentially finite-type extension, the chosen
diagonal element maps to the base selector \(q_H\).  Under the stated
perfectness and rank hypotheses, the normal-closure base change is an
equivalence \(N\otimes_KL\simeq (G/H\to N)\).  This is a generic statement, not
a compactified extension.

\begin{leanfiles}
These comparisons are in \normsrc{SelectorComparison.lean}, especially
\leanpath{NormalClosureData.collisionToGenericSheets_elem_eq_baseSheetSelector}
and \leanpath{NormalClosureData.normalClosureTensorGenericSheetsEquiv}.
\end{leanfiles}

\section{Compactification: raw defect and a sufficient graph route}
\label{sec:compactification}

Let \(\Geta\) and \(\Gbar\) be generic and compactified rational
correspondence algebras, and let \(r:\Gbar\to\Geta\) be an algebra map.  For
an idempotent \(e_\eta\in\Geta\), a raw closure package has a witness type
\(\Cl_r(e_\eta)\) and assigns to each \(x\in\Cl_r(e_\eta)\) a class
\(\Pi_x\in\Gbar\) satisfying \(r(\Pi_x)=e_\eta\).  Its defect is

\[
  \Def(\Pi_x)=\Pi_x^2-\Pi_x.
\]

\begin{proposition}[Raw closure defect --- \formal; \conditional]
\label{prop:raw-defect}
Every raw closure satisfies
\[
  r(\Def(\Pi_x))=0.
\]
Moreover \(\Def(\Pi_x)=0\) if and only if \(\Pi_x\) is idempotent.  The
zero-defect subtype \(\Cl_r^0(e_\eta)\) therefore gives an actual extension
package \(\Ext_r(e_\eta)\).
\end{proposition}

Kernel membership is the entire unconditional conclusion.  No current
theorem identifies \(\ker r\) with boundary-supported Chow classes or kills
the particular defect.

\begin{statusnote}{\sourceinput/\openstatus}
In the dimension-theoretic setup of the cited result, let \(X\) be locally
of finite type over the base \(S\), and let \(j:U\hookrightarrow X\) be an
open immersion with reduced complement \(i:Z\hookrightarrow X\).  Then the
Chow localization sequence
\[
  \CH_k(Z)\xrightarrow{i_*}\CH_k(X)
    \xrightarrow{j^*}\CH_k(U)\longrightarrow0
\]
is exact \cite{StacksChowLocalization}.  This supplies the standard
existence of a raw Chow-class lift and identifies the ambiguity between lifts
as boundary-supported.  It does not say that the nonlinear class
\(\Def(\Pi_x)\) vanishes or that a chosen lift respects correspondence
composition.
\end{statusnote}

A stronger supplied datum gives a different route.  Suppose generic and
compactified spaces carry actions of the same group and that their free graph
algebras map to \(\Geta\) and \(\Gbar\).  A compactified deck-graph
comparison is the square

\[
\begin{CD}
  \Q[G] @>{\gamma_{\mathrm{cpt}}}>> \Gbar\\
  @V{\gamma_\eta}VV @VV{r}V\\
  \Geta @= \Geta,
\end{CD}
\]

where the two arrows out of \(\Q[G]\) are the supplied graph-cycle maps.  We
denote the comparison datum by \(\mathscr G_r\); its compatibility equation is
\(r\gamma_{\mathrm{cpt}}=\gamma_\eta\).

\begin{proposition}[Compactified graph image --- \formal; \conditional]
\label{prop:graph-route}
Given \(\mathscr G_r\), every idempotent \(e\in\Q[G]\) has a compactified
image
\(\Pi=\gamma_{\mathrm{cpt}}(e)\) satisfying
\[
  \Pi^2=\Pi,
  \qquad r(\Pi)=\gamma_\eta(e),
  \qquad \Def(\Pi)=0.
\]
The subgroup average \(e_H\) yields the corresponding singleton-witness
extension package.
\end{proposition}

For smooth projective schemes over a fixed field \(k\), composition of rational Chow
correspondences is associative and bilinear
\cite{StacksChowCorrespondenceComposition}, while graph classes compose
contravariantly \cite{StacksChowGraphFunctor}.  Once an actual compactified
deck action and its restriction comparison have been constructed, these
standard results supply the Chow-theoretic multiplication laws used by the
graph route, after matching that variance convention to the project's graph
algebra map.  They do not construct the action, compactification, comparison,
or convention match.

\begin{leanfiles}
Raw closures and Proposition~\ref{prop:raw-defect} are in
\boundarysrc{ProjectorClosureDefect.lean}, under
\leanpath{AlgebraProjectorClosure}, \leanpath{convolutionDefect}, and
\leanpath{toAlgebraProjectorExtension}.  The supplied graph square and
Proposition~\ref{prop:graph-route} are in
\boundarysrc{CompactifiedGraphProjector.lean}, under
\leanpath{CompactifiedDeckGraphComparison} and
\leanpath{groupAlgebraProjectorExtension}.
\end{leanfiles}

\begin{statusnote}{\openstatus}
Neither a compactification carrying the required deck action nor a concrete
Chow-valued graph square \(\mathscr G_r\) is constructed.  The sufficient
route explains exactly why global deck graphs would bypass a separate
boundary-defect vanishing proof; it does not assert that such graphs exist.
\end{statusnote}

\section{Residual realization, coverage, and the exact gap}
\label{sec:coverage}

Fix the action \(\rho:\Gbar\to\End_{\Q}(\Q\oplus V)\).  Supply an extension
package for \(q_H\), choose its compactified projector \(\Pi\), and supply a
linear residual realization \(\theta:E_H\to V\).  Put
\(\Theta=\operatorname{id}_{\Q}\oplus\theta\) and
\(\alpha=\theta(a_H)\).  The comparison square required by the active route
is

\begin{equation}
\label{eq:action-square}
  \rho(\Pi)\Theta=\Theta P_{a_H}.
\end{equation}

\begin{proposition}[Marked-selector realization --- \formal; \conditional]
\label{prop:realization}
From Equation~\eqref{eq:action-square},
\[
  d_{\mathrm{ll}}(\rho(\Pi))=\alpha.
\]
If \(\theta\) is surjective, then \(\rho(\Pi)=P_\alpha\).  If the ambient
action also has the Chow lower-left detection map of
Section~\ref{sec:target}, then \(\alpha\) lies in the cycle-class image.
\end{proposition}

For a fixed cycle-class model and ambient action, define
\(\operatorname{Cov}_{\mathrm{Gal}}(\mathscr H,\rho)\) to mean that every
rational Hodge class \(\alpha\) has some class-dependent marked-selector
witness.  Its finite group \(G\), subgroup \(H\), extension package,
compactified projector, residual map, and comparison may all depend on
\(\alpha\), but \(\Gbar\) and \(\rho\) remain fixed.

\begin{theorem}[Coverage implications --- \formal; \conditional]
\label{thm:coverage}
A detection map and
\(\operatorname{Cov}_{\mathrm{Gal}}(\mathscr H,\rho)\) imply
\(\mathrm{HC}_{\Q}(\mathscr H)\).  Conversely, construction and detection
give the graph-projector equivalence of
Theorem~\ref{thm:hodge-graph-equivalence}; to replace an arbitrary graph lift
by marked Galois data and obtain
\[
  \mathrm{HC}_{\Q}(\mathscr H)
  \Longleftrightarrow
  \operatorname{Cov}_{\mathrm{Gal}}(\mathscr H,\rho),
\]
one must additionally assume the explicitly named strong replacement
hypothesis.
\end{theorem}

\begin{leanfiles}
The comparison record and its realization consequences are in
\chowsrc{GaloisMarkedRealization.lean}, under
\leanpath{MarkedGaloisSelectorComparison},
\leanpath{lowerLeftClass_action_eq}, and
\leanpath{isAlgebraicHodgeClass_of_comparison}.  Class-dependent coverage,
the one-way Hodge implication, the strong replacement premise, and the
conditional equivalence are respectively
\leanpath{ClassDependentMarkedGaloisSelectorCoverage},
\leanpath{rationalHodgeConjectureAt_of_classDependentMarkedGaloisSelectorCoverage},
\leanpath{StrongGraphProjectorLiftReplacementByClassDependentMarkedGaloisSelector},
and \leanpath{rationalHodgeConjectureAt_iff_classDependentMarkedGaloisSelectorCoverage}.
\end{leanfiles}

\subsection*{Exact remaining gap --- \openstatus}

To obtain an actual Hodge-theoretic result by this route, the following
constructions must be supplied for genuine geometry, with their displayed
compatibilities proved rather than postulated:

\begin{enumerate}[leftmargin=2.2em]
  \item genuine Chow groups, Betti cohomology, Hodge subspaces, and the
        rational cycle-class map, together with the Tate-block construction,
        lower-left Chow detection, and Hodge-morphism comparison;
  \item for every targeted Hodge class, marked normal-closure data and a
        geometric residual realization \(\theta:E_H\to V\) carrying
        \(a_H\) to that class;
  \item either a compactified deck-graph square \(\mathscr G_r\), or a raw
        closure plus a proof that its specific defect vanishes;
  \item the action-intertwining identity
        \(\rho(\Pi)\Theta=\Theta P_{a_H}\), and enough surjectivity when the
        entire graph operator rather than only its lower-left class is needed;
  \item a proof of the class-dependent coverage assertion.  For a literal
        equivalence with coverage, one also needs the strong replacement of
        arbitrary graph lifts by marked Galois witnesses.
\end{enumerate}

The codimension-one theorem of Section~\ref{sec:boundary} can feed such a
construction only after a genuine global purity or gluing theorem is supplied.
It does not close items (2)--(5) by itself.

\section{Lean theorem and module map}
\label{sec:lean-map}

\renewcommand{\arraystretch}{1.2}
\begingroup
\small
\begin{longtable}{@{}>{\raggedright\arraybackslash}p{0.25\textwidth}
                >{\raggedright\arraybackslash}p{0.70\textwidth}@{}}
\textbf{Mathematical layer} & \textbf{Principal declaration and Lean module}\\ \hline
\endfirsthead
\textbf{Mathematical layer} & \textbf{Principal declaration and Lean module}\\ \hline
\endhead
Rational Hodge target (\formal) &
\leanpath{RationalHodgeConjectureAt}\newline
\leanpath{Realization/HodgeClasses.lean}\\
Graph-projector equivalence (\formal; \conditional) &
\leanpath{rationalHodgeConjectureAt_iff_everyHodgeClassHasGraphProjectorLift}\newline
\leanpath{Realization/Chow/TateGraphProjector.lean}\\
Collision decomposition (\formal) &
\leanpath{affineSelfCollisionDecompositionOfFormallyUnramified}\newline
\leanpath{Collision/AffineGeometry.lean}\\
Free finite-\'etale algebra (\formal) &
\leanpath{FiniteEtaleFreeAlgebra}\newline
\leanpath{Correspondence/EtaleClasses.lean}\\
Group projector transport (\formal) &
\leanpath{groupAlgebraCorrespondence_isIdempotent}\newline
\leanpath{Correspondence/GroupAlgebraProjector.lean}\\
Hecke projector transport (\formal) &
\leanpath{heckeCorrespondence_isIdempotent}\newline
\leanpath{Correspondence/HeckeProjector.lean}\\
Marked block normal form (\formal) &
\leanpath{decomposition_conj_multiplication}\newline
\leanpath{Boundary/Sheets/MarkedBlockProjector.lean}\\
Shared block core (\formal) &
\leanpath{operator_eq_shear_diagonal_shear_symm}\newline
\leanpath{LinearAlgebra/BlockGraphProjector.lean}\\
Conjugate-boundary endpoint (\formal; \conditional) &
\leanpath{noCodimensionOneRamification_of_conjugateBoundarySeparation}\newline
\leanpath{Boundary/Normalization/ConjugateBoundary.lean}\\
Subgroup extension (\formal; \conditional) &
\leanpath{CompactifiedSubgroupProjectorExtension}\newline
\leanpath{Boundary/GaloisProjectorExtension.lean}\\
Raw closure defect (\formal; \conditional) &
\leanpath{restriction_convolutionDefect}\newline
\leanpath{Boundary/ProjectorClosureDefect.lean}\\
Compactified graph route (\formal; \conditional) &
\leanpath{groupAlgebraProjectorExtension}\newline
\leanpath{Boundary/CompactifiedGraphProjector.lean}\\
Marked residual comparison (\formal; \conditional) &
\leanpath{MarkedGaloisSelectorComparison}\newline
\leanpath{Realization/Chow/GaloisMarkedRealization.lean}\\
Coverage equivalence (\formal; \conditional) &
\leanpath{rationalHodgeConjectureAt_iff_classDependentMarkedGaloisSelectorCoverage}\newline
\leanpath{Realization/Chow/GaloisMarkedRealization.lean}\\
\end{longtable}
\endgroup

All displayed paths in the table are relative to the
\leanpath{HodgeConjecture/} directory.  Umbrella modules collect these
declarations but do not add mathematical hypotheses.

\appendix

\section{Technical appendix: localization, fixed rings, and inertia}
\label{app:localization}

This appendix records the local commutative-algebra bridge used when a
normalization diagram is expressed with valuation inertia.  It is not the
main projector construction.

Let a finite group \(G\) act on a domain \(S\), let \(I\leq H\leq G\), and
fix a height-one prime \(Q\subset S\).  With
\(q_J=Q\cap S^J\), contraction gives

\[
  (S^G)_{q_G}\longrightarrow(S^H)_{q_H}
    \longrightarrow(S^I)_{q_I}.
\]

Finite invariant extensions are integral; the formalization proves the
required prime transitivity and going-down consequences, preservation of
height one, the DVR property after localizing a Noetherian normal fixed ring,
and the canonical local injection from the \(H\)-fixed to the \(I\)-fixed
local ring.  Formal unramifiedness descends along this selected tower under
the explicit source premise.

For \(S=\operatorname{IntCl}_B(N)\), let \(I_Q^{\mathrm{pr}}\) be the kernel
of the stabilizer action on the residue ring at \(Q\).  If a centered
discrete valuation ring \(V\subset N\) is identified with \(S_Q\), then its
valuation inertia \(I_V\) equals \(I_Q^{\mathrm{pr}}\).  The equality uses
the localization comparison; the two inertia groups are not definitionally
the same object.

\begin{statusnote}{\sourceinput/\conditional}
In the standard finite-group invariant setting, the Stacks Project
fixed-prime-inertia lemma states that the map from the full invariant ring to
the inertia-fixed ring is \'etale at the selected prime
\cite{StacksPrimeInertiaEtale}; the DVR specialization is recorded separately
in \cite{StacksDVRInertiaEtale}.  The Lean endpoint takes the corresponding
localized formal-unramifiedness statement as an explicit premise because the
external theorem is not in the pinned Mathlib API.  We use the cited theorem
as a source input and do not reproduce its proof here.
\end{statusnote}

\begin{leanfiles}
The local fixed-ring descent is in \normsrc{FixedRingLocalization.lean},
\normsrc{FixedRingUnramified.lean}, and
\normsrc{PrimeInertiaDescent.lean}.  Prime inertia is in
\normsrc{PrimeInertia.lean}.  The valuation localization and equality of
inertia groups are in \normsrc{LocalizationBridge.lean}, principally
\leanpath{primeInertiaComparisonAt_of_isNoetherianRing} and
\leanpath{valuationInertia_eq_normalizationPrimeInertiaGroup_of_isNoetherianRing}.
\end{leanfiles}

\section{Source inputs and notation audit}
\label{app:sources}

The machine-checked development uses Lean~4 and Mathlib
\cite{Lean4,Mathlib}.  Standard commutative-algebra background used by the
local appendix includes integral-closure normality
\cite{StacksIntegralClosureNormal}, localization of Noetherian and normal
rings \cite{StacksNoetherianLocalization,StacksNormalLocalization}, and the
one-dimensional local criterion for a DVR \cite{StacksDVRCriterion}.  The
fixed-field/integral-closure comparison may be compared with the normal
Galois integral-closure statement \cite{StacksGaloisIntegralClosure}; that
citation is context and does not itself prove the project-specific algebra
equivalence.  These citations supply standard mathematical context; the
corresponding algebraic lemmas used in the principal spine are proved in Lean
or obtained from Mathlib.

\begin{statusnote}{\formal}
For a finite-type ring map, the cited Stacks Project result detects the
non-unramified locus by containment of the Noether different
\cite{StacksNoetherDifferent}.  Lean proves the collision-presentation
statement under the localization-friendly, essentially finite-type
hypothesis as
\leanpath{collisionDifferentDetectsRamificationAt_of_essFiniteType}; its
proof is not repeated in this report.  The citation verifies the ordinary
finite-type subcase and is not an unencoded premise of the Lean theorem.
This detection theorem does not
construct inertia comparisons, a compactification, a purity map, a Chow
class, or a prescribed cohomological action.
\end{statusnote}

The following non-identifications govern all notation in this manuscript:

\[
\begin{array}{rcl}
  q_H\in F_{\Q}(H) &\neq& e_H\in\Q[G],\\
  e_H\in\Q[G] &\neq& \Pi\in\Gbar,\\
  \partial U\subset Y_L &\neq& \partial_g\subset Y_N,\\
  I_V &\neq& I_Q^{\mathrm{pr}}\quad\text{definitionally},\\
  \Cl_r(e_\eta) &\neq& \Ext_r(e_\eta),\\
  \operatorname{Ram}^{(1)}(Y_N/B)=\varnothing
    &\neq& \text{global selector-gluing purity}.
\end{array}
\]

Each equality or implication between two rows requires the comparison map or
hypothesis displayed in the relevant theorem.

\section{Conclusion}

The formal development now has an explicit endpoint and an auditable gap.
The rational Hodge target is exactly equivalent to canonical graph-projector
lifting once construction and lower-left detection are supplied.  Collision
algebra, finite-\'etale convolution, group and Hecke idempotents, the marked
sheet block calculation, and the conjugate-boundary codimension-one theorem
provide compatible algebraic and geometric layers toward such lifts.  Raw
compactified closure is deliberately weaker than projector extension; a
supplied compactified deck-graph square is a clean sufficient alternative.

What remains is geometric rather than terminological: instantiate the
genuine Chow/Hodge model, construct class-dependent residual realizations and
compactified projectors, prove the action square and detection compatibility,
and establish coverage.  Until those inputs are constructed, the repository
is a conditional formalization architecture and does not prove a Hodge case.

\paragraph{Acknowledgment.}
OpenAI Codex assisted with the Lean formalization and manuscript draft.  The
manuscript remains subject to author review; responsibility for its final
statements, proofs, citations, and release terms belongs to the author.

\begingroup
\small
\bibliographystyle{alpha}
\bibliography{references}
\endgroup

\end{document}
